% !TEX root = ../b-rg.tex
\chapter{Adjectives}\label{ch:adjectives}

Borlish adjectives have two forms: a base form and a comparative form. The comparative is formed by adding a suffix, which shifts stress to the penultimate syllable.

\section{Formation of the comparative}\label{sec:comparative-formation}

\subsection{Regular comparatives in \bl{-essem}}
The regular comparative suffix is \bl{-essem}, which replaces or follows the final segment of the adjective and attracts penultimate stress.
    \glphm{leger legeressem}{leˈʒɛr ˌle.ʒeˈrɛ.sɛm}{light lighter}
Several orthographic adjustments may occur when the suffix is added. Before the vowel-initial suffix, \orth{ç} becomes \orth{c}, while \orth{c} becomes \orth{ch} to preserve the \phm{k} pronunciation.
    \glphm{lavaç lavacessem | starc starchessem}{laˈvats ˌla.vaˈtsɛ.sɛm | stark starˈkɛ.sɛm}{wasteful {more wasteful} | strong stronger}
After a final orthographic \orth{-q}, a linking \orth{-u-} is inserted.
    \glphm{picq picquessem}{pɪk pɪˈkwɛ.sɛm}{small smaller}
The final consonant may be doubled in the comparative form to maintain the checked quality of the preceding vowel. This occurs most commonly with final \orth{-l}.
    \glphm{boncop boncoppessem}{bɔnˈkɔp ˌbɔn.kɔˈpɛ.sɛm}{cheap cheaper}
The final consonant may be voiced in the comparative form. This is seen most commonly in present participles ending in \orth{-nt} (which form comparatives in \orth{-ndessem}) and in adjectives ending in fricatives, though only final \orth{-f} reflects this orthographically as \orth{-vessem}.
    \glphm{tart tardessem | braf bravessem}{tart tarˈdɛ.sɛm | braf braˈvɛ.sɛm}{late later | brave braver}
Adjectives ending in a written \orth{-Vu} diphthong form the comparative in \orth{-Vvessem}.
    \glphm{blau blavessem}{blo blaˈvɛ.sɛm}{blue bluer}
Certain historically long vowels in the final syllable of an adjective may shorten in the comparative form. This occurs most commonly with adjectives ending in \bl{-ous}.
    \glphm{fatuous fatuosessem}{faˈtjuz ˌfa.tjoˈzɛ.sɛm}{clumsy clumsier}
Adjectives ending in \orth{-e} form the comparative in \orth{-eyessem}.
    \glphm{idone idoneyessem}{ˌi.doˈne iˌdo.niˈjɛ.sɛm}{suitable {more suitable}}

\subsection{Adjectives in syllabic \bl{-r}}
Adjectives ending in syllabic \orth{-r} show several patterns. Occasionally they form regular comparatives in \bl{-ressem}.
    \glphm{jundr jundressem}{ˈʒɪn.dr̩ ʒɪnˈdrɛ.sɛm}{junior {more junior}}
Some adjectives, most notably those ending in the suffix \bl{-abr} `-able', form comparatives in \bl{-lessem} rather than \bl{-ressem}.
    \glphm{stabr stablessem | sengr senglessem}{ˈstab.r̩ staˈblɛ.sɛm | ˈsɛn.gr̩ sɛnˈglɛ.sɛm}{steady steadier | simple simpler}
The majority of adjectives ending in syllabic \orth{-r} form the comparative with the suffix \bl{-errem}, which replaces the final syllabic consonant and is stressed penultimately. The same adjustments that affect \bl{-essem} may also affect \bl{-errem}.
    \glphm{nuvr nuverrem | diçr dicerrem}{ˈnɪv.r̩ niˈvɛ.rɛm | ˈdɪts.r̩ diˈtsɛ.rɛm}{fresh fresher | great greater}
    \glphm{allagr allagherrem | leur leverrem}{aˈlɛj.r̩ ˌa.lɛˈjɛ.rɛm | ˈlau.r̩ leˈvɛ.rɛm}{happy happier | free freer}

\subsection{Irregular comparatives}
A few adjectives show irregular stem changes in the comparative.
    \glphm{brey bressem}{bri ˈbrɛ.sɛm}{short shorter}
Several common adjectives have fully suppletive comparative forms.
    \glphm{bon ottem | mal pessem | gran maissem}{bɔn ˈɔ.tɛm | mal ˈpɛ.sɛm | gran ˈme.sɛm}{good better | bad worse | big bigger}

% --- Planned sections ---
% \section{Position of the adjective}
%   Mostly postnominal (y rout novel, stal zajadau); which adjectives
%   precede (bon? gran? picq?) and any meaning differences.
% \section{Attributive and predicative use}
%   Any differences in form or position; choice of copula with
%   predicative adjectives (es vs sta: permanent vs temporary?).
% \section{Adjective complements and modifiers}
%   Adjectives taking arguments (capabr de ...); degree adverbs before or
%   after the adjective.
% (Comparative and superlative constructions: see the Comparison chapter.)
% \section{Participles as adjectives}
%   Past participles in -að (and irregulars), present participles in
%   -ant/-ent; their comparatives (-ndessem, see above).
% \section{Adjectives used as nouns and adverbs}
%   Substantivised adjectives; bare adjective as adverb (Scaunç glosses
%   'tap-pst def drum agile and exact', 'wash-pst def face quick'); cos + adjective (cos nuvr 'lately', cos unic, cos equal).
%   Cross-reference the Adverbs chapter.
